A transformer language model is iterated in three ways: over layers, where the residual stream of each token is updated step by step; over optimizer steps, where the weights evolve under gradient descent; and over generated tokens, where the model feeds its own samples back as input. We call these the three clocks of a language model. Each clock has its own literature on clustering, edge of stability, loss spikes, grokking, repetition loops and temperature transitions, and these literatures rarely cite each other or the nonlinear dynamics they draw on. We place the three clocks in one framework of coupled maps, review about 200 papers, and condense the results into the Three-Clock Dictionary, which maps fifteen concepts from nonlinear dynamics onto each clock and marks every correspondence as proved, measured or conjectured. Experiments on open models, all run on a laptop CPU, put these tools to work. Training makes each layer of a Pythia model stretch perturbations more strongly, while the forward pass as a whole changes little. At the edge of stability of a small transformer, two to six Hessian eigenvalues exceed $2/\eta$ while the curvature along the step stays at $2/\eta$. Escape from repetition loops in GPT-2 and Pythia follows the Arrhenius law of Kramers' escape problem, with the logit gap as the barrier (correlation 0.997 between fitted and measured barriers), and the barrier deepens while the model repeats itself and then saturates. A loop therefore never locks for good, but the temperature needed to leave it within a budget of $N$ cycles falls like $1/\ln N$, and it rises with model size, by about half from 70M to 1B parameters. Coherent text occupies a narrow window of temperatures between repetition and incoherence. A small reasoning model, Qwen3-0.6B on grade-school math, shows the same window, with the best accuracy at the temperature its developers recommend, but at low temperature it rethinks the same step again and again without entering an exact loop. We close with twelve open problems and release code, data and a maintained reading list.
